\chapter{Introduction}\label{cha: introduction}
Recently, there has been an interest in developing \gls{THz} \gls{PCA} that can be activated by lasers operating at telecommunication wavelengths of \qty{1.55}{\micro\meter} due to the possibility of integration with existing technology\cite{baconPhotoconductiveEmittersPulsed2021}. These \gls{PCA}s are used in \gls{THz} \gls{TDS} for various applications, including biology\cite{yuMedicalApplicationTerahertz2019}, imaging\cite{petrovApplicationTerahertzPulse2016}, and pharmaceuticals\cite{dohiApplicationTerahertzPulse2016}, among others. A \gls{PCA} consists of a semiconductor material where the terahertz generation or absorption occurs and two metal electrodes on the semiconductor surface separated by a few microns that forms the photoconductive gap\cite{songHandbookTerahertzTechnologies2015}. When a femtosecond pulse from the laser illuminates the gap, carriers are generated in the semiconductor and swept by an electric field from the bias applied to the metal electrodes, and the movement of the carriers generates the terahertz wave. 

One of the common semiconductor materials of choice for a \gls{PCA} operating at the \qty{1.55}{\micro\meter} wavelength is \gls{LT} grown indium gallium arsenide (InGaAs), which has a bandgap of $E_{g}<$ \qty{0.8}{\electronvolt} and can be grown lattice matched to an indium phosphide (InP) substrate via molecular beam epitaxy or \gls{MBE}\cite{burfordReviewTerahertzPhotoconductive2017}. However, when an InGaAs PCA is illuminated by a laser there are losses involved due to the high refractive index of the material\cite{riturajHighlySensitiveEfficient2025}. Adding an antireflective coating to the InGaAs PCA can increase the efficiency of the device by reducing the amount of light reflected by the film. The challenge lies in the design of a suitable antireflective coating for InGaAs.

\section{Significance of the study}\label{sec: significance}
The development of \gls{THz} \gls{PCA}s compatible with telecommunication wavelengths holds great potential for integrating \gls{THz} technology with existing optical systems. By enhancing the efficiency of \gls{THz} generation, this study contributes to the advancement of \gls{THz} \gls{TDS} systems, which have diverse applications in fields such as biology, imaging, and pharmaceuticals. The current use of InP substrates for \gls{MBE} growth of \gls{LT}-InGaAs based \gls{PCA}s allows creation of films operating at \qty{1.55}{\micro\meter} wavelengths with acceptable strain. However, this setup faces limitations, including \gls{THz} absorption by the InP substrate and inefficient laser energy absorption by the InGaAs thin film. By integrating a \gls{DBR} stack, the study aims to address these problems, improving the \gls{THz} intensity generated by the \gls{PCA} and its overall performance. The findings could pave the way for more efficient \gls{THz} devices, enhancing capabilities in scientific research, medical diagnostics, and industrial processes.

\section{Objectives of the study}\label{sec: Objectives}
The main objective of this study is to investigate the feasibility of using a-Si and SiO$_2$ thin films to create an antireflective coating which can be used for \gls{LT}-InGaAs \gls{PCA}s. Specifically, this will require the following objectives:
\begin{itemize}
	\item to identify the optical properties of a-Si and SiO$_2$ thin films deposited using electron beam evaporation,
	\item to design an antireflective coating utilizing a-Si and SiO$_2$ optimized for the wavelengths at which \gls{LT}-InGaAs THz \gls{PCA}s can be operated, and
	\item to evaluate the reflectivity of the fabricated a-Si/SiO$_2$ antireflective coating at the relevant wavelengths.
\end{itemize}

\section{Scope and limitations of the study}\label{sec: sco}

The study will only cover the investigation of an \gls{a-Si} and SiO$_2$ \gls{ARC} for \gls{LT}-InGaAs THz \gls{PCA} applications. As such, the study will not investigate the growth of the \gls{LT}-InGaAs  film using \gls{MBE} and will also be limited to electron beam evaporation as the deposition method for the \gls{a-Si}/SiO$_2$. The characterization is also limited to the evaluation of the reflectivity of the designed films and will not include applications of the developed system.